% GENERATED FILE. Edit canonical lessons, then run npm run generate:books.
% canonical-source-hash: fnv1a64:f85a971e06f2ba54
% canonical-lessons: ZH-C238-tiankong, ZH-C238-yueliang, ZH-C238-wuyun, ZH-C238-dahai, ZH-C238-caihong

\chapter{Sea and Sky}
\label{ch:zh-sea-and-sky}
\hlchaptermodality{238}

\begin{chapteropening}
\textbf{By the end of this chapter:} \emph{I can describe the sky, the moon, dark clouds, the open sea and a rainbow.}

Complete the last lesson of chapter 238: I can describe the sky, the moon, dark clouds, the open sea and a rainbow.
\end{chapteropening}

\section[tiānkōng]{tiānkōng — the sky}

\label{lesson:ZH-C238-tiankong}



\begin{quote}
Before the new one: say the Mandarin for to sprout, then the Mandarin for a watermelon.
\end{quote}

\subsection*{You'll want to know: tiānkōng}
\emph{tiānkōng} — \textquotedblleft{}the sky\textquotedblright{}.

\begin{quote}
\emph{jīntiān de tiānkōng hěn lán} — The sky is so blue today.
\end{quote}

Its characters come later; for now it is a word for the ear.

\begin{cousinweb}
\emph{tiānkōng} is built from \zh{天} 'heaven' and \emph{kōng} 'empty', both inherited from Middle Chinese, \emph{then} and \emph{khuwng} (Guangyun, in Baxter's transcription), which glosses the second 'hollow'.
\end{cousinweb}

\begin{grammarlens}[title={What you've built}]
One more word for describing things.
\end{grammarlens}

\subsection*{Your turn}
\begin{itemize}
\raggedright
  \item \emph{Say it:} \emph{tiānkōng}
  \item \emph{Say it:} \emph{tiānkōng}, once more
  \item \emph{Say it:} \emph{jīntiān de tiānkōng hěn lán} — The sky is so blue today.
  \item \emph{Recall:} say the Mandarin for to sprout, then the Mandarin for a watermelon, then say \emph{tiānkōng} again
\end{itemize}

\begin{tcolorbox}[breakable,colback=teal!4,colframe=teal!35!black,title={Before you move on}]
What does tiānkōng mean? (\textquotedblleft{}The sky\textquotedblright{}.) Say it once more.
\end{tcolorbox}
\section[yuèliang]{yuèliang — the moon}

\label{lesson:ZH-C238-yueliang}



\begin{quote}
Before the new one: say the Mandarin for a watermelon, then the Mandarin for the sky.
\end{quote}

\subsection*{You'll want to know: yuèliang}
\emph{yuèliang} — \textquotedblleft{}the moon\textquotedblright{}.

\begin{quote}
\emph{jīntiān wǎnshang yuèliang hěn dà} — The moon is huge tonight.
\end{quote}

Its characters come later; for now it is a word for the ear.

\begin{cousinweb}
\emph{yuèliang} is built from \zh{月} 'moon' and \emph{liang} 'bright', both inherited from Middle Chinese, \emph{ngjwot} and \emph{ljangH} (Guangyun, in Baxter's transcription), which glosses the second 'bright'; in the word the second loses its tone.
\end{cousinweb}

\begin{grammarlens}[title={What you've built}]
One more word for describing things.
\end{grammarlens}

\subsection*{Your turn}
\begin{itemize}
\raggedright
  \item \emph{Say it:} \emph{yuèliang}
  \item \emph{Say it:} \emph{yuèliang}, once more
  \item \emph{Say it:} \emph{jīntiān wǎnshang yuèliang hěn dà} — The moon is huge tonight.
  \item \emph{Recall:} say the Mandarin for a watermelon, then the Mandarin for the sky, then say \emph{yuèliang} again
\end{itemize}

\begin{tcolorbox}[breakable,colback=teal!4,colframe=teal!35!black,title={Before you move on}]
What does yuèliang mean? (\textquotedblleft{}The moon\textquotedblright{}.) Say it once more.
\end{tcolorbox}
\section[wūyún]{wūyún — dark clouds}

\label{lesson:ZH-C238-wuyun}



\begin{quote}
Before the new one: say the Mandarin for the sky, then the Mandarin for the moon.
\end{quote}

\subsection*{You'll want to know: wūyún}
\emph{wūyún} — \textquotedblleft{}dark clouds\textquotedblright{}.

\begin{quote}
\emph{wūyún lái le yào xiàyǔ le} — Dark clouds are coming in; it's about to rain.
\end{quote}

Its characters come later; for now it is a word for the ear.

\begin{cousinweb}
\emph{wūyún} is built from \emph{wū} 'black as a crow' and \zh{云} 'cloud'. The second is inherited from Middle Chinese \emph{hjun} (Guangyun, in Baxter's transcription); the dictionary explains the first, \emph{'u} there, after an older book, as the crow that is all black.
\end{cousinweb}

\begin{grammarlens}[title={What you've built}]
One more word for describing things.
\end{grammarlens}

\subsection*{Your turn}
\begin{itemize}
\raggedright
  \item \emph{Say it:} \emph{wūyún}
  \item \emph{Say it:} \emph{wūyún}, once more
  \item \emph{Say it:} \emph{wūyún lái le yào xiàyǔ le} — Dark clouds are coming in; it's about to rain.
  \item \emph{Recall:} say the Mandarin for the sky, then the Mandarin for the moon, then say \emph{wūyún} again
\end{itemize}

\begin{tcolorbox}[breakable,colback=teal!4,colframe=teal!35!black,title={Before you move on}]
What does wūyún mean? (\textquotedblleft{}Dark clouds\textquotedblright{}.) Say it once more.
\end{tcolorbox}
\section[dàhǎi]{dàhǎi — the open sea}

\label{lesson:ZH-C238-dahai}



\begin{quote}
Before the new one: say the Mandarin for the moon, then the Mandarin for dark clouds.
\end{quote}

\subsection*{You'll want to know: dàhǎi}
\emph{dàhǎi} — \textquotedblleft{}the open sea\textquotedblright{}.

\begin{quote}
\emph{wǒ hěn xǐhuan dàhǎi} — I really love the sea.
\end{quote}

Its characters come later; for now it is a word for the ear.

\begin{cousinweb}
\emph{dàhǎi} is built from \zh{大} 'big' and \emph{hǎi} 'sea', both inherited from Middle Chinese, \emph{daH} and \emph{xojX} (Guangyun, in Baxter's transcription).
\end{cousinweb}

\begin{grammarlens}[title={What you've built}]
One more word for describing things.
\end{grammarlens}

\subsection*{Your turn}
\begin{itemize}
\raggedright
  \item \emph{Say it:} \emph{dàhǎi}
  \item \emph{Say it:} \emph{dàhǎi}, once more
  \item \emph{Say it:} \emph{wǒ hěn xǐhuan dàhǎi} — I really love the sea.
  \item \emph{Recall:} say the Mandarin for the moon, then the Mandarin for dark clouds, then say \emph{dàhǎi} again
\end{itemize}

\begin{tcolorbox}[breakable,colback=teal!4,colframe=teal!35!black,title={Before you move on}]
What does dàhǎi mean? (\textquotedblleft{}The open sea\textquotedblright{}.) Say it once more.
\end{tcolorbox}
\section[cǎihóng]{cǎihóng — rainbow}

\label{lesson:ZH-C238-caihong}



\begin{quote}
Before the new one: say the Mandarin for dark clouds, then the Mandarin for the open sea.
\end{quote}

\subsection*{You'll want to know: cǎihóng}
\emph{cǎihóng} — \textquotedblleft{}rainbow\textquotedblright{}.

\begin{quote}
\emph{yǔ tíng le yǒu cǎihóng} — The rain's stopped, and there's a rainbow.
\end{quote}

Its characters come later; for now it is a word for the ear.

\begin{cousinweb}
\emph{cǎihóng} is built from \emph{cǎi} 'colour' and \emph{hóng} 'rainbow'. The first is inherited from Middle Chinese \emph{tshojX} (Guangyun, in Baxter's transcription), glossed 'brilliance'; the dictionary glosses the second, \emph{huwng} there, with an old two-character word for the rainbow.
\end{cousinweb}

\begin{grammarlens}[title={What you've built}]
One more word for describing things.
\end{grammarlens}

\subsection*{Your turn}
\begin{itemize}
\raggedright
  \item \emph{Say it:} \emph{cǎihóng}
  \item \emph{Say it:} \emph{cǎihóng}, once more
  \item \emph{Say it:} \emph{yǔ tíng le yǒu cǎihóng} — The rain's stopped, and there's a rainbow.
  \item \emph{Recall:} say the Mandarin for dark clouds, then the Mandarin for the open sea, then say \emph{cǎihóng} again
\end{itemize}

\begin{tcolorbox}[breakable,colback=teal!4,colframe=teal!35!black,title={Before you move on}]
What does cǎihóng mean? (\textquotedblleft{}Rainbow\textquotedblright{}.) Say it once more.
\end{tcolorbox}
